\documentclass[11pt]{article}
\usepackage[utf8]{inputenc}
\usepackage[spanish]{babel}
\usepackage{amsmath,amssymb,amsthm}
\usepackage{graphicx}
\usepackage[margin=2.5cm]{geometry}
\usepackage{enumitem}
\title{Semana 6 -- Ejercicios}
\author{Tutorías -- Introducción al Cálculo}
\date{07/09--11/09/2026}
\begin{document}
\maketitle

\section{Ejercicios de la semana}

\begin{enumerate}
    \item Determine si $F = \{(x,y) \in \mathbb{R}\times\mathbb{R} : y^2 = x\}$ es función.
    Justifique.

    \item Sea $F = \{(x,y)\in\mathbb{R}\times\mathbb{R} : y = 3x-2\}$. Determine si $F$ es función
    de $\mathbb{R}$ en $\mathbb{R}$ y calcule $\text{DOM}\, F$ y $\text{REC}\, F$.

    \item Sea $A = \{1,2,3\}$ y $F = \{(1,2),(2,2),(3,5)\}$. ¿Es $F$ función de $A$ en
    $\mathbb{R}$? Determine $\text{REC}\, F$.

    \item Determine si $F$ y $G$ son la misma función, donde $F(x) = x^2-1$ con
    $\text{DOM}\, F = \mathbb{R}$ y $G(x) = (x-1)(x+1)$ con $\text{DOM}\, G = \mathbb{R}$.

    \item Grafique $F(x) = 3x-4$, indicando la intersección con cada eje.

    \item Grafique $F(x) = -x^2+2x+3$, indicando el vértice de la parábola.

    \item Grafique $f(x) = |x-2|$.

    \item Determine el mayor subconjunto de $\mathbb{R}$ en el que
    $F(x) = \sqrt{4-x^2}$ define una función (es decir, el dominio máximo posible).
    \hfill (Resp.\ $[-2,2]$)

    \item \textbf{(Nivel interrogación)} Sea $F(x) = \sqrt{x^2-9}$. Determine el dominio máximo
    de $F$, grafique $F$ indicando de qué curva conocida es un tramo, y determine su recorrido.

    \item \textbf{(Nivel interrogación)} Sean $F(x) = \sqrt{9-x^2}$ y $G(x) = -\sqrt{9-x^2}$,
    ambas con dominio $[-3,3]$. Explique, usando la definición de función, por qué $F$ y $G$ son
    funciones por separado, pero la relación $F\cup G = \{(x,y) : x\in[-3,3],\ y=F(x)\ \text{o}\
    y=G(x)\}$ no lo es. ¿Qué representa geométricamente $F\cup G$?

    \item \textbf{(Nivel interrogación)} Determine si
    $H = \{(x,y)\in\mathbb{R}\times\mathbb{R} : y=x^2 \text{ si } x\ge0,\ y=-x^2 \text{ si }
    x<0\}$ es función. En caso afirmativo, determine su dominio y recorrido, y grafique.
\end{enumerate}

\end{document}
